% ============================================================
% 分片：s05  对应原文 PDF 第 66--68 页（自第 66 页中的附录 B 标题起）
% 负责章节：附录 B "$A$-type Quantum Mechanics and Morse Theory" + B.1 + B.2 + B.3
% 说明：\appendix 已在 s04 中给出，本分片的 \section 自动编号为 B。
% 译者：translator-06   审校状态：初稿
% ============================================================

% ==== 原文 PDF 第 66 页（本页上半为附录 A 的收尾，属 s04；本分片自附录 B 标题起） ====
\section{$A$ 型量子力学与 Morse 理论}

\begin{keytakeaway}
\textbf{核心物理图像：}一维 $\mathcal{N}=2$（$A$ 型）量子力学的玻色势能为 $\sim|\partial h|^2$，故基态集中在超势 $h$ 的临界点附近；当 $h$ 为 Morse 函数时，每个临界点对应唯一的微扰基态，理论的基态数即 $h$ 的临界点数。
\textbf{数学工具：}$\mathcal{N}=2$ 超空间 $\mathbb{R}^{1|2}$ 上的超场与超导数、按超空间积分构造作用量、超势的 Hessian、Morse 函数、单一超荷的 $\bar Q$-上同调。
\textbf{主要结论：}基态即同时 $Q$-闭与 $\bar Q$-闭的零能态（BPS 态），可仅由一个超荷的上同调 $H(\mathcal{H},\bar Q)=\operatorname{Ker}\bar Q/\operatorname{Im}\bar Q$ 完整刻画；这为把基态计数化为 Morse 理论提供了场论框架。
\end{keytakeaway}

本附录回顾一维 $\mathcal{N}=2$ 理论（称为 $A$ 型量子力学）的一些基本特征，以及它与 Morse 理论的关系。如 2.4.3 节所述，该理论是在二维 $\mathcal{N}=(2,2)$ 理论的 $A$ 扭变中研究拓扑 $A$ 膜（topological $A$-brane）的出发点之一；事实上，[48] 的分析正是借助这一视角的教益来处理整个 $A$ 扭变二维理论的代数性质，例如编码体与边界局域算子碰撞的 $L A_\infty$ 结构。本附录的材料改编自 [9, Section 10] 与 [48, Section 10]；我们也推荐把 $A$ 型量子力学与 Morse 理论联系起来的原始工作 [47]。
% [存疑: 原文 "encoding colliding of bulk and boundary local operators" 中 "colliding" 疑为名词 "collisions" 之误，译文按名词理解。]

\subsection{$\mathcal{N}=2$ 超空间}

用 $\mathcal{N}=2$ 超空间 $\mathbb{R}^{1|2}$ 描述 $\mathcal{N}=2$ 量子力学理论是一种方便的语言。我们在 Euclidean $\mathbb{R}^{1|2}$ 上引入坐标 $t,\theta,\bar\theta$。对 $t$ 的导数记作 $\partial_t$，对 $\theta$ 与 $\bar\theta$ 的导数记作 $\partial_\theta$ 与 $\partial_{\bar\theta}$。它们满足
\begin{equation}\tag{B.1}
\{\partial_\theta,\theta\} = \{\partial_{\bar\theta},\bar\theta\} = [\partial_t,t] = 1,
\end{equation}
% [Tutor Note: 易错预警：$\partial_\theta,\partial_{\bar\theta}$ 与坐标的（反）对易关系把 $\theta,\bar\theta$ 视为幂零的奇坐标，而 $t$ 与 $\partial_t$ 仍是通常的偶变量；由此任何超场沿 $\theta,\bar\theta$ 的 Taylor 展开都只有有限项。写分量作用量时须保持这一反交换约定，否则 (B.6) 第二行的符号会出错。]
其余（反）对易子皆为零。我们采用如下约定：Grassmann 奇变量的复共轭交换次序，$(\chi\xi)^*=\bar\xi\bar\chi$。

超空间平移的生成元由哈密顿量 $H=-\partial_t$\footnote{在 Lorentz 号差下，我们把哈密顿量与矢量场 $H=-i\partial_{x^0}$ 等同；Wick 转动到 Euclidean 号差 $t=ix^0$ 意味着 $H=-\partial_t$。}与超荷
\begin{equation}\tag{B.2}
Q := \partial_\theta + \bar\theta\partial_t \qquad \bar Q := -\partial_{\bar\theta} - \theta\partial_t.
\end{equation}
给出。我们还需要手征与反手征超导数
\begin{equation}\tag{B.3}
D := \partial_\theta - \bar\theta\partial_t \qquad \bar D := -\partial_{\bar\theta} + \theta\partial_t.
\end{equation}

% ==== 原文 PDF 第 67 页 ====
这些算子合起来满足如下关系：
\begin{equation}\tag{B.4}
\{\bar Q,Q\} = -2\partial_t = 2H \qquad \{\bar D,D\} = 2\partial_t = -2H,
\end{equation}
其余（反）对易子皆为零。存在一个 $U(1)_R$ $R$ 对称性，它作用在 $\theta$ 上荷为 $-1$、作用在 $\bar\theta$ 上次数为 1；这意味着 $Q$ 的 $R$ 荷为 1，$\bar Q$ 的 $R$ 荷为 $-1$。
% [存疑: 原文此处对 $\bar\theta$ 用 "degree 1"，而前文与后文均用 "charge"（荷），疑为用词不一致；译文照原样保留“次数”。]

超空间上的函数，或超空间上更一般丛的截面，称为超场。一般超场 $S$ 可按奇变量 $\theta,\bar\theta$ 作 Taylor 展开；由于这些奇变量幂零，该展开是有限和：
\begin{equation}\tag{B.5}
S = a + \theta\alpha - \bar\theta\bar\beta + \theta\bar\theta b.
\end{equation}
此外，超场可以是玻色子的，即 $a,b$ 为玻色子而 $\alpha,\beta$ 为费米子；也可以是费米子的，即 $a,b$ 为费米子而 $\alpha,\beta$ 为玻色子。

$\mathcal{N}=2$ 超对称代数通过矢量场 $Q,\bar Q$ 作用于超场，并在分量场上诱导出作用。例如有
\begin{equation}\tag{B.6}
\begin{aligned}
QS &= \alpha + \bar\theta(b+\partial_t a) + \theta\bar\theta(-\partial_t\alpha)\\
&:= Qa - \theta Q\alpha + \bar\theta Q\bar\beta + \theta\bar\theta Qb
\end{aligned}
\end{equation}
于是 $Qa=\alpha$、$Q\alpha=0$、$Q\bar\beta=b+\partial_t a$、$Qb=-\partial_t\alpha$。第二行中额外出现的符号来自把费米子对称性 $Q$ 移过 $\theta$ 与 $\bar\theta$。
% [Tutor Note: 推导补足：第二行只是把第一行按 $\theta,\bar\theta$ 的各次项重新归类，并把系数写成 $Q$ 作用于相应分量场的形式。例如 $\bar\theta(b+\partial_t a)$ 对应 $\bar\theta Q\bar\beta$，而 $\theta\bar\theta(-\partial_t\alpha)$ 对应 $\theta\bar\theta Qb$；符号差异全部来自 $Q$（费米子）与 $\theta,\bar\theta$ 的反交换。]

一般超场 $S$ 给出超对称代数的可约表示，可以对 $S$ 施加条件以得到不可约超场。下文将用到满足 $X^*=X$ 的（玻色子）实超场 $X$：
\begin{equation}\tag{B.7}
X = x + \theta\xi - \bar\theta\bar\xi + \theta\bar\theta f.
\end{equation}
$\mathcal{N}=2$ 超对称代数在这些场上的作用因此为
\begin{equation}\tag{B.8}
\begin{aligned}
Qx &= \xi &\qquad Q\xi &= 0 &\qquad Q\bar\xi &= \partial_t x + f &\qquad Qf &= \partial_t\bar\xi\\
\bar Qx &= \bar\xi &\qquad \bar Q\xi &= \partial_t x - f &\qquad \bar Q\bar\xi &= 0 &\qquad \bar Qf &= -\partial_t\xi
\end{aligned}
\end{equation}
自然地给超场 $X$ 赋 $U(1)_R$ 荷 0，于是 $x,f$ 的荷为 0，$\xi$ 的荷为 1，$\bar\xi$ 的荷为 $-1$。我们最终关心的是上同调分次 $C=-R$，使得扭变超荷 $\bar Q$ 具有上同调次数 1。
% [Tutor Note: 物理图像：取 $C=-R$ 使 $\bar Q$ 成为次数 $+1$ 的扭变超荷，正是 $A$ 型扭变在一维的对应物；此后 $\bar Q$-上同调就是理论的物理态空间（基态），而 $Q$ 则扮演“另一半”超对称的角色。]

\subsection{$A$ 型量子力学}

把超场在超空间上积分，可以写下明显超对称的作用量。设映射 $X$ 取值于带度规 $g_{ab}(X)$ 的黎曼流形 $M$；我们的 $\mathcal{N}=2$ $\sigma$ 模型的自然动能项如下：
\begin{equation}\tag{B.9}
\begin{aligned}
S_{kin} &= \frac{1}{2}\int dt\,d^2\theta\left( g_{ab}(X)DX^a\bar DX^b \right)\\
&= \int dt\left( \frac{1}{2}g_{ab}\partial_t x^a\partial_t x^b - g_{ab}\bar\xi^a\partial_t\xi^b - \frac{1}{2}R_{abcd}\bar\xi^a\xi^b\bar\xi^c\xi^d \right.\\
&\qquad\qquad\qquad \left. - \frac{1}{2}g_{ab}\left( f^a+\Gamma^a{}_{cd}\bar\xi^c\xi^d \right)\left( f^b+\Gamma^b{}_{ef}\bar\xi^e\xi^f \right) \right)
\end{aligned}
\end{equation}
% [Tutor Note: 推导补足：这里用了超空间积分的约定 $\int d^2\theta\,\theta\bar\theta=1$（其余项积分为零）。于是 $g_{ab}(X)DX^a\bar DX^b$ 中 $\theta\bar\theta$ 的系数给出 $-g_{ab}\bar\xi^a\partial_t\xi^b$，展开 $g_{ab}(X)$ 到 $\theta\bar\theta$ 阶并用 Riemann 张量重整化即产生曲率项 $-\frac{1}{2}R_{abcd}\bar\xi^a\xi^b\bar\xi^c\xi^d$；辅助场 $f$ 只出现在最后一项。]

% ==== 原文 PDF 第 68 页 ====
由于附近唯一自然的手征超场是 $\bar DX$，我们实质上只能考虑来自在整个超空间上积分的相互作用。选取光滑实函数 $h:M\to\mathbb{R}$（也称为超势），并在超空间上积分它
\begin{equation}\tag{B.10}
S_h = -\int dt\,d^2\theta\, h(X) = \int dt\left( -f^a\partial_a h(x) + \partial_a\partial_b h(x)\bar\xi^a\xi^b \right)
\end{equation}
由此得到完整作用量为这两项之和
\begin{equation}\tag{B.11}
S = S_{kin} + S_h = \int dt\,d^2\theta\left( \frac{1}{2}g_{ab}(X)DX^a\bar DX^b - h(X) \right).
\end{equation}
由此可见 $f^a$ 是辅助场，其运动方程把它们取定为 $f^a=-g^{ab}\partial_b h-\Gamma^a{}_{bc}\bar\xi^b\xi^c$。重要的是，消去 $f$ 后的玻色势能为 $\sim|\partial h|^2=g^{ab}(\partial_a h)(\partial_b h)$，因此基态将局域在 $h$ 的临界点附近，即满足对每个 $a$ 都有 $\partial_a h=0$ 的位置。在这样的临界点 $p$ 的邻域内，费米子 $\bar\xi^a,\xi^a$ 的质量矩阵依赖于 $h$ 的 Hessian，$\sim\partial_a\partial_b h(p)$。我们将限制在任意临界点邻域内都没有无质量费米子的超势 $h$\footnote{这一限制意味着，在微扰论最低阶，每个临界点会有一个唯一的基态。若存在无质量费米子，则费米子真空 $|0\rangle$（被例如 $\xi$ 湮灭）与 $\bar\xi|0\rangle$ 都是零能态。}，即 Hessian $\sim\partial_a\partial_b h(p)$ 在每个临界点 $p$ 处非退化。这样的 $h$ 称为 Morse 函数。
% [存疑: 原文 "the fermions ... have a mass matrix depends on the Hessian of h" 语法有误，应为 "whose mass matrix depends on..."；译文按后者理解。]
% [Tutor Note: 物理图像：消去辅助场 $f$ 后，作用量中只留下势能项 $\sim|\partial h|^2$，故能量最低的场构型集中在 $h$ 的临界点；$h$ 为 Morse 函数（Hessian 非退化）时，各临界点附近的质量谱非零且互不混合，因而可以有意义地逐点计数基态。]
% [背景补充: Morse 函数要求所有临界点非退化（Hessian 可逆）；此时每个临界点可赋予一个指标（Hessian 负特征值的个数），这正是 Morse 理论中构造 Morse 复形、并以梯度流连接临界点的数据。]

\subsection{BPS 态与上同调}

在量子化上述理论之前，我们先对 $\mathcal{N}=2$ 理论中 Hilbert 空间的一些结构性质作评论。我们的目标是辨别这类理论中基态的独特之处。非常一般地，设存在费米子对称性 $Q$、$\bar Q=Q^\dagger$，即 $[H,Q]=[H,\bar Q]=0$，且 $[Q,\bar Q]=Q\bar Q+\bar Q Q=2H$，其中 $[-,-]$ 为分次对易子，$H$ 为量子理论的哈密顿量。一个直接的结构性要点是：由幺正性界，态的能量必然非负：设 $|\psi\rangle$ 为单位范数的能量 $E$ 本征态，$\langle\psi|\psi\rangle=1$，则
\begin{equation}\tag{B.12}
E = E\langle\psi|\psi\rangle = \langle\psi|H|\psi\rangle = \frac{1}{2}\langle\psi|[Q,\bar Q]|\psi\rangle = \frac{1}{2}\langle Q\psi|Q\psi\rangle + \frac{1}{2}\langle\bar Q\psi|\bar Q\psi\rangle \ge 0,
\end{equation}
其中 $|Q\psi\rangle=Q|\psi\rangle$、$|\bar Q\psi\rangle=\bar Q|\psi\rangle$。特别地，态 $|\psi\rangle$ 是基态当且仅当其能量为零 $E=0$，而这当且仅当 $Q|\psi\rangle=0$ 且 $\bar Q|\psi\rangle=0$。若 $Q|\psi\rangle=0$ 则称 $|\psi\rangle$ 为 $Q$-闭，类似地若 $\bar Q|\psi\rangle=0$ 则称其为 $\bar Q$-闭；因此基态同时是 $Q$-闭与 $\bar Q$-闭的。这些态常称为 BPS 态。
% [Tutor Note: 推导补足：最后的不等号只需 $H=\frac{1}{2}[Q,\bar Q]$ 的矩阵元表示 $E=\frac{1}{2}\left(\|Q\psi\|^2+\|\bar Q\psi\|^2\right)$；这是“超对称保证能谱非负”的普遍机制，与 $\mathcal{N}=2$ 的具体动力学无关。]

结果表明，我们可以仅用一个超荷（例如 $\bar Q$）及相应的 $\bar Q$-闭态来编码基态。注意 $\bar Q^2=\frac{1}{2}[\bar Q,\bar Q]=0$ 意味着：若 $|\psi\rangle$ 是 $\bar Q$-闭的，则对任意 $\psi'$，$|\psi\rangle+\bar Q|\psi'\rangle$ 也是 $\bar Q$-闭的。形如 $\bar Q|\psi'\rangle$ 的态称为 $\bar Q$-恰当的；$\bar Q$-恰当态必然 $\bar Q$-闭。我们于是构造 $\bar Q$-上同调 $H(\mathcal{H},\bar Q)$ 为 $\bar Q$-闭态模去 $\bar Q$-恰当态所成的空间：
\begin{equation}\tag{B.13}
H(\mathcal{H},\bar Q) := \operatorname{Ker}\bar Q/\operatorname{Im}\bar Q = \{|\psi\rangle\in\mathcal{H}\,|\,\bar Q|\psi\rangle=0\}/\left(|\psi\rangle\sim|\psi\rangle+\bar Q|\psi'\rangle\right).
\end{equation}
% [Tutor Note: 物理图像：$H(\mathcal{H},\bar Q)$ 的每个上同调类恰对应一个基态——$\bar Q$-闭给出零能条件，而模去 $\bar Q$-恰当态正是去掉零模中的纯规范方向；这正是把基态计数化为 Morse 理论的关键一步。]
